\section{Lie--Poisson reduction and canonical realization}
\label{app:lie}

\subsection{Graded Lie algebra}

Let
\[
\mathfrak n
=
\operatorname{span}
\{P,Q,K,V,U,Y,X\}
\]
with nonzero brackets
\[
[P,Q]=K,
\]
\[
[P,K]=V,\qquad [Q,K]=U,
\]
\[
[P,U]=X,\qquad [P,V]=Y,
\]
\[
[Q,U]=-Y,\qquad [Q,V]=X.
\]

The grading is
\[
\mathfrak n_1=\operatorname{span}\{P,Q\},\quad
\mathfrak n_2=\operatorname{span}\{K\},
\]
\[
\mathfrak n_3=\operatorname{span}\{V,U\},\quad
\mathfrak n_4=\operatorname{span}\{Y,X\},
\]
so the growth vector is
\[
\boxed{(2,3,5,7).}
\]

Equivalently,
\[
[P,[P,[P,Q]]]+[Q,[Q,[P,Q]]]=0.
\]

\subsection{Lie--Poisson coordinates}

Let
\[
(h_1,h_2,\kappa,v,u,\eta,\xi)
\]
be dual to
\[
(P,Q,K,V,U,Y,X).
\]

Using
\[
\{\ell_A,\ell_B\}=\ell_{[A,B]},
\]
the nonzero brackets are
\[
\{h_1,h_2\}=\kappa,
\]
\[
\{h_1,\kappa\}=v,\qquad
\{h_2,\kappa\}=u,
\]
\[
\{h_1,u\}=\xi,\qquad
\{h_1,v\}=\eta,
\]
\[
\{h_2,u\}=-\eta,\qquad
\{h_2,v\}=\xi.
\]

The central coordinates \(\xi,\eta\) are Casimirs. A third Casimir is
\begin{equation}
\boxed{
\mathcal C
=
\kappa(\xi^2+\eta^2)
-\frac12(v^2-u^2)\eta
-vu\xi.
}
\label{eq:casimir}
\end{equation}

Thus generic coadjoint orbits are four-dimensional.

\subsection{Rotational reduction and scaling}

Let
\[
r=\sqrt{\xi^2+\eta^2}>0.
\]
Using the \(SO(2)\) automorphism, choose
\[
\xi=r,\qquad \eta=0.
\]
Then
\[
\mathcal C=\kappa r^2-vur,
\]
so
\[
\kappa=\frac{\mathcal C}{r^2}+\frac{uv}{r}.
\]

Under Carnot dilation,
\[
h_i\mapsto\lambda h_i,\qquad
\kappa\mapsto\lambda^2\kappa,
\]
\[
u,v\mapsto\lambda^3(u,v),\qquad
\xi,\eta\mapsto\lambda^4(\xi,\eta).
\]
Since
\[
\mathcal C\mapsto\lambda^{10}\mathcal C,
\]
choosing \(\lambda=r^{-1/4}\) normalizes the central momentum. The dimensionless parameter is
\[
\boxed{
\mu=\frac{\mathcal C}{r^{5/2}}.
}
\]
Dropping scaled hats,
\[
\xi=1,\qquad \eta=0,\qquad
\boxed{
\kappa=\mu+uv.
}
\]

\subsection{Canonical realization}

Let
\[
(u,v,\Pi_u,\Pi_v)
\]
be canonical coordinates with
\[
\{u,\Pi_u\}=1,\qquad
\{v,\Pi_v\}=1.
\]

The realization compatible with the Lie--Poisson convention is
\[
\boxed{
h_1=-\Pi_u,
}
\]
\[
\boxed{
h_2=-\Pi_v+\mu u+\frac12u^2v.
}
\]

The normal Hamiltonian becomes
\[
\boxed{
H_\mu
=
\frac12
\left[
\Pi_u^2+
\left(
\Pi_v-\mu u-\frac12u^2v
\right)^2
\right].
}
\]

Define
\[
w=\Pi_v-\mu u-\frac12u^2v.
\]
Then
\[
\dot u=\Pi_u,\qquad
\dot v=w,
\]
\[
\dot\Pi_u=(\mu+uv)w,\qquad
\dot\Pi_v=\frac12u^2w.
\]
Therefore
\[
\boxed{
\ddot u=(\mu+uv)\dot v,\qquad
\ddot v=-(\mu+uv)\dot u.
}
\]

\subsection{Unit-speed reduction}

On \(H_\mu=\tfrac12\), write
\[
\dot u=\cos\phi,\qquad
\dot v=\sin\phi.
\]
Then
\[
\boxed{
\dot\phi=-(\mu+uv).
}
\]

Introduce
\[
x=\frac{u+v}{\sqrt2},\qquad
y=\frac{v-u}{\sqrt2},\qquad
z=\phi-\frac{\pi}{4}.
\]
Since
\[
uv=\frac{x^2-y^2}{2},
\]
the reduced system is
\[
\boxed{
\dot x=\cos z,\qquad
\dot y=\sin z,\qquad
\dot z=-\frac{x^2}{2}+\frac{y^2}{2}-\mu.
}
\]
